\documentclass[11pt,a4paper]{article}

\usepackage[margin=1in]{geometry}
\usepackage[T1]{fontenc}
\usepackage{lmodern}
\usepackage{booktabs}
\usepackage{array}
\usepackage{amsmath}
\usepackage{xcolor}
\usepackage{listings}
\usepackage{enumitem}
\usepackage[hidelinks]{hyperref}
\usepackage{titlesec}

\titleformat{\section}{\normalfont\large\bfseries}{\thesection}{0.6em}{}
\titleformat{\subsection}{\normalfont\normalsize\bfseries}{\thesubsection}{0.6em}{}
\setlist[itemize]{leftmargin=1.4em,itemsep=2pt,topsep=3pt}
\setlist[enumerate]{leftmargin=1.6em,itemsep=2pt,topsep=3pt}

\definecolor{codebg}{gray}{0.96}
\lstset{
  basicstyle=\ttfamily\footnotesize,
  backgroundcolor=\color{codebg},
  frame=single,
  framerule=0pt,
  breaklines=true,
  columns=fullflexible,
  keepspaces=true,
  xleftmargin=0.5em,
  showstringspaces=false,
  aboveskip=0.7em,
  belowskip=0.7em
}

\newcommand{\code}[1]{\texttt{\small #1}}

\title{\vspace{-2em}\textbf{CALM: Calibration-Aware Logits Mixing for Efficient}\\
\textbf{and Trustworthy Medical Image Segmentation}\\[0.4em]
\large Research Proposal --- Path A (Lab Paper Extension / Methodological Enhancement)}
\author{Md. Irtiza Hossain \\
\small M.S. Computer Science \& Engineering, BRAC University \\
\small \href{mailto:mohammad.irtiza.hossain@gmail.com}{mohammad.irtiza.hossain@gmail.com}}
\date{6 September 2026}

\begin{document}
\maketitle

\noindent\begin{tabular}{@{}ll@{}}
\textbf{Builds on} & LoMix (NeurIPS 2025), EMCAD (CVPR 2024), G-CASCADE (WACV 2024) \\
\textbf{Preliminary results} & From my own replication of LoMix on Synapse (Deliverables 1--2) \\
\textbf{Harness} & \url{https://github.com/irtiza1999/LoMix-Reproducibility-Study} \\
\end{tabular}

\section{Motivation and significance}

Segmentation models are increasingly proposed for point-of-care use --- rural
clinics, bedside ultrasound, endoscopy suites --- where the hardware is modest
and no specialist is standing by to catch a bad contour. In that setting a model
needs two properties at once. It must be \textbf{cheap at inference}, and it
must \textbf{know when it is unsure}, because an overconfident wrong boundary on
an organ at risk is worse than an abstention.

LoMix is a notably elegant answer to the first requirement. By supervising
\emph{mixtures} of multi-scale decoder logits under learned weights, it improves
segmentation while adding \textbf{nothing} to the inference graph: the mixing
weights are parameters of the loss module, not the network. I verified this
directly during replication --- all four supervision arms share a byte-identical
inference graph (26.773\,M parameters, 4.324\,GMACs, 36.8\,ms at $224^2$ on a
GTX 960).

But LoMix answers the second requirement not at all, and its treatment of the
first is incomplete in a specific way. The learned mixing weights are
\textbf{global and input-independent}: training produces one fixed answer to
``which scale do I trust'', applied identically to a clean abdominal slice and
to an ambiguous pancreas boundary under scanner shift. Clinically, those are not
the same question. Methodologically, a supervision scheme that already reasons
about \emph{which} of its own predictions to believe is one conditioning step
away from reasoning about \emph{when} to believe them --- which is the
calibration question.

This proposal takes that step. \textbf{CALM} makes the mixing weights a function
of per-sample uncertainty, derived from signals the decoder already computes,
and adds a calibration objective alongside Dice. It keeps LoMix's defining
property --- the inference graph is untouched --- while addressing three
concrete weaknesses I measured in the released implementation.

\section{Problem definition and gaps}

\subsection{Gap 1 --- supervision weights are static and input-independent}

LoMix learns one scalar per (subset, operation) pair and freezes it after
training. Under domain shift, or on the ambiguous boundaries where segmentation
actually fails, the useful scale mixture plausibly differs per sample. Nothing
in the formulation can express that.

My replication supplies evidence that this is not merely a theoretical concern.
Starting from a common initialisation of $\mathrm{softplus}(0) = 0.6931$, after
20 epochs the four per-scale weights had drifted to 0.379, 0.386, 0.392, 0.400
--- a spread of 0.021 across scales, against a movement of ${\sim}0.3$ away from
initialisation shared by all of them. The weights moved \emph{together}, not
\emph{apart}. Whatever differentiation the mechanism eventually achieves, at
this budget it had barely begun, which suggests one global scalar per subset is
a very slow-learning parameterisation for the quantity it is trying to capture.

\subsection{Gap 2 --- combinatorial cost falls entirely on training}

For $K$ decoder outputs and $|O|$ operations, LoMix evaluates
$(2^K - K - 1)\times|O|$ mixed maps --- 44 for the released $K=4$, $|O|=4$
configuration --- each with its own CE and Dice term. Measured on identical
hardware and batch:

\begin{center}
\begin{tabular}{@{}lrrr@{}}
\toprule
\textbf{Arm} & \textbf{Peak train VRAM} & \textbf{Step time} &
\textbf{vs single-logit} \\
\midrule
\code{last\_layer} & 1008\,MiB & 649\,ms & $1.00\times$ \\
\code{deep\_supervision} & 1054\,MiB & 708\,ms & $1.09\times$ \\
\code{mutation} (unweighted mixing) & 1268\,MiB & 938\,ms & $1.45\times$ \\
\code{lomix} (learned weights) & 2290\,MiB & 1649\,ms & $\mathbf{2.54\times}$ \\
\bottomrule
\end{tabular}
\end{center}

The paper reports inference cost, which is genuinely zero. It does not report
this. On constrained hardware, $2.54\times$ training cost decides whether the
method can be used at all --- and it scales as $2^K$, so deeper decoders make it
worse.

\subsection{Gap 3 --- the objective is accuracy-only}

LoMix optimises CE + Dice. Neither is a calibration objective, and the paper
reports no calibration metric. A method whose entire premise is \emph{weighting
predictions by how much they should be trusted} never measures whether the
resulting confidence is trustworthy.

\subsection{Gap 4 --- two of the four mixing operations are defective as released}

Found and verified during replication (\code{tools/diagnose\_loss\_ops.py}).

\textbf{\code{mul} is numerically unbounded.} It multiplies raw logits, so a
$k$-map combination is a $k$-fold product growing like $|\text{logit}|^k$. On
the trained checkpoint's own logits ($\max|p| = 6.16$):

\begin{center}
\begin{tabular}{@{}lrrr@{}}
\toprule
\textbf{Operation} & $k=2$ & $k=3$ & $k=4$ \\
\midrule
\code{add} & 12.0 & 17.9 & 23.6 \\
\textbf{\code{mul}} & \textbf{36.2} & \textbf{211.0} & \textbf{1214.0} \\
\code{wf} & 6.1 & 6.1 & 5.9 \\
\code{concat} & 2.8 & 2.9 & 1.9 \\
\bottomrule
\end{tabular}
\end{center}

In my run this diverged to NaN at epoch 13 under the paper's own
hyperparameters, destroying the arm.

\textbf{\code{concat} cannot learn.} Its $1\times1$ fusion convolution is
constructed \emph{inside} \code{forward}, randomly re-initialised on every call,
never registered as a submodule and never given to the optimizer. Two identical
forward passes under \code{no\_grad} differ by 5.41 for \code{concat} and by
exactly 0.0 for \code{add}, \code{mul} and \code{wf}. Consistent with this,
\code{concat}'s eleven learned weights collapsed to a common 0.3755--0.3770 and
stopped moving --- there is no gradient signal for them to differentiate on.

So of four mixing operations, one injects unbounded values and one injects
random noise. Any claim about \emph{which} mixtures matter rests on an operation
set that is half broken.

\subsection{Related work and what is missing}

\textbf{Deep supervision and multi-scale fusion} (UNet++, CASCADE, EMCAD,
G-CASCADE, LoMix) weight scales by fixed or globally-learned coefficients. None
conditions on the input.

\textbf{Uncertainty estimation} (MC dropout, deep ensembles, evidential
segmentation) produces per-sample uncertainty but at multiplied inference cost
--- the opposite of what deployment on a point-of-care device needs.

\textbf{Calibration} (temperature scaling, focal loss, proper scoring rules) is
applied post-hoc or as a loss term, but is not connected to \emph{how
supervision is allocated across scales}.

\textbf{Conditional computation / routing} (mixture-of-experts, dynamic
networks) conditions computation on the input, but in the network, so inference
cost rises.

The gap CALM occupies: \textbf{per-sample conditioning of the supervision, not
the network.} Because the conditioning lives entirely in the loss, it costs
nothing at inference --- which no uncertainty-aware or conditional-computation
method above can say.

\section{Proposed method}

\subsection{Overview}

\begin{lstlisting}
                                    +----------- training only ------------+
 encoder -- decoder --+-- p4 --+    |                                      |
                      |-- p3 --|    |  u(x) = [entropy, cross-scale JSD,   |
                      |-- p2 --|----+->        boundary density]           |
                      +-- p1 --+    |             |                        |
                          |         |             v                        |
                          |         |     g_theta(u)   small MLP           |
                          |         |             |                        |
                          |         |             v                        |
                          |         |  per-sample weights w(x)_{S,o}       |
                          |         |             |                        |
                          |         |  Gumbel top-m subset selection       |
                          |         |             |                        |
                          |         |             v                        |
                          |         |  sum of w(x) * [CE + Dice + calib]   |
                          |         +--------------------------------------+
                          |
 inference:  encoder -- decoder -- p1     <- unchanged, zero added cost
\end{lstlisting}

\subsection{Bounded mixing operations (prerequisite)}

Before conditioning anything, the operation set must be sound.

\begin{itemize}
\item \textbf{\code{mul} in probability space.} Multiply per-class softmax
probabilities and renormalise (equivalently, \emph{average logits} --- a
log-domain product), which is a principled product-of-experts fusion and is
bounded by construction. This removes the divergence mode in \S2.4 at no
representational cost.
\item \textbf{\code{concat} registered properly.} Move the $1\times1$ fusion
convolutions into an \code{nn.ModuleDict} keyed by subset size, as the \code{wf}
operation already correctly does. This makes a quarter of the operation set
functional for the first time.
\end{itemize}

Both changes carry falsifiable predictions: \code{mul} should stop diverging,
and \code{concat}'s learned weights should move away from their collapsed value.
If \code{concat} remains at the collapsed value once trainable, that is itself
an informative result about which mixtures matter.

\subsection{A free per-sample uncertainty signal}

The key design constraint is that uncertainty must not cost an extra forward
pass, or the efficiency argument collapses. CALM therefore derives $u(x)$ from
quantities the decoder has already produced:

\begin{enumerate}
\item \textbf{Predictive entropy} of the final map, mean over foreground pixels.
\item \textbf{Cross-scale disagreement} --- mean Jensen--Shannon divergence
between the softmax outputs of the four decoder scales. This is exactly the
signal that should govern how much to trust each scale, and it is available for
free because all four maps are already computed for the loss.
\item \textbf{Boundary density} --- a cheap proxy for the ambiguous-boundary
regime, from the gradient magnitude of the predicted mask.
\end{enumerate}

Cross-scale disagreement is the conceptually important one: it makes the
weighting responsive to \emph{the specific failure mode the mixing is meant to
fix}.

\subsection{Conditioning the weights}

Replace LoMix's global scalar $w_{S,o}$ with
\[
  w(x)_{S,o} \;=\; \mathrm{softplus}\!\left( a_{S,o} + b_{S,o}^{\top}\,
  g_\theta\big(u(x)\big) \right)
\]
where $a_{S,o}$ is the original global parameter (so CALM strictly generalises
LoMix, recovering it exactly when $b = 0$), $g_\theta$ is a small MLP
($3 \to 16 \to 8$), and $b_{S,o}$ is a learned per-term projection. Total added
parameters ${\approx}10^3$, all in the loss module, none at inference.
Initialising $b = 0$ means training starts exactly at LoMix and departs only if
conditioning helps --- which makes the comparison against the baseline clean
rather than confounded by initialisation.

\subsection{Uncertainty-guided subset selection}

Rather than evaluating all 44 mixed maps every step, sample $m \ll 44$ subsets
per step via Gumbel top-$m$ over the conditioned weights, with a
straight-through estimator. Terms the model considers uninformative for
\emph{this} sample are skipped. This attacks \S2.2 directly and is the mechanism
by which CALM aims to be \emph{cheaper} than LoMix in training, not merely
better. Expected cost at $m = 12$ is roughly 1.4--1.6$\times$ the single-logit
baseline versus LoMix's measured $2.54\times$, to be confirmed empirically.

\subsection{Calibration in the objective}

Add a differentiable calibration term to each supervised map --- a soft-binned
ECE surrogate or a proper scoring rule --- weighted by $\lambda$. Report ECE and
reliability diagrams alongside Dice. This turns ``which predictions to trust''
from an implicit assumption into a measured quantity.

\section{Experimental plan}

\textbf{Datasets.} Synapse multi-organ and ACDC (in-distribution, directly
comparable to the lab's published tables); polyp datasets (Kvasir-SEG,
CVC-ClinicDB, CVC-ColonDB, ETIS) for a \textbf{cross-centre domain-shift}
protocol --- train on two centres, test on the held-out ones, which is where
per-sample conditioning should matter most and where a static global weight
should not.

\textbf{Baselines.} \code{last\_layer}, \code{deep\_supervision},
\code{mutation}, \code{lomix} (as released), \code{lomix + bounded ops}, and
CALM ablations: bounded-ops only, +conditioning, +selection, +calibration term.

\textbf{Metrics.} Dice and HD95; \textbf{ECE, adaptive ECE, reliability
diagrams, AURC}; and training cost (step time, peak memory) measured under
identical conditions, since a cost claim is part of the contribution.

\textbf{Key ablations.} (i) Which uncertainty feature carries the signal ---
entropy alone vs cross-scale JSD alone vs all three. (ii) Sensitivity to $m$.
(iii) Does conditioning help more under shift than in-distribution --- the
central prediction. (iv) Do the conditioned weights vary meaningfully across
samples, or collapse back to a global solution? A negative here would falsify
the premise, and I would report it.

\textbf{Compute.} In-distribution Synapse/ACDC runs are ${\sim}7$\,h per arm at
reduced budget on the hardware I have used so far; the full protocol needs
cluster access, which I would expect to plan around rather than assume.

\section{Preliminary results}

All from my own LoMix replication (Deliverable 1); every figure traces to a run
log, and the logs, the patched source and the analysis tools are public at
\url{https://github.com/irtiza1999/LoMix-Reproducibility-Study}. Reduced budget:
20 of the repository's default 300 epochs, effective batch 6, seed 2222.

\begin{center}
\begin{tabular}{@{}lrl@{}}
\toprule
\textbf{Arm} & \textbf{Best mean Dice} & \textbf{Status} \\
\midrule
\code{last\_layer} & 0.6764 & completed \\
\code{deep\_supervision} & \textbf{0.7979} & completed \\
\code{mutation} & 0.7788 & completed \\
\code{lomix} & 0.7225 (pre-divergence) & diverged (NaN) at epoch 13 \\
\bottomrule
\end{tabular}
\end{center}

For reference, the paper's Table~1 reports the same four arms on the same
backbone as 80.9 / 82.9 / 83.6 / 85.1 DICE, averaged over at least three runs on
an RTX A6000.

\textbf{These runs do not reproduce the paper's ordering, and I do not claim
they refute it.} At 20 of the default 300 epochs, with one seed, the mixing
weights had
barely differentiated (\S2.1) --- the mechanism that is supposed to produce
LoMix's advantage had not yet engaged. What the runs \emph{do} establish,
independently of budget, is: the zero-inference-overhead claim (measured
directly), the $2.54\times$ training-cost ratio (measured under identical
conditions), and the two implementation defects in \S2.4 (verified on the
trained model's own logits, not dependent on training length).

Those three facts are precisely the premises this proposal builds on, which is
why the replication was worth doing before proposing anything.

\section{Expected outcomes, risks and timeline}

\textbf{Expected outcomes.} (1) A supervision scheme that conditions on
per-sample uncertainty at zero inference cost and \emph{lower} training cost
than LoMix. (2) Improved calibration under domain shift at equal or better Dice.
(3) A corrected, properly-registered operation set, with an empirical account of
which mixtures actually matter once all four operations work.

\textbf{Risks and mitigations.}
\begin{itemize}
\item \emph{Uncertainty estimates are themselves miscalibrated early in
training.} Mitigate with a warm-up during which $b = 0$ (pure LoMix), enabling
conditioning only once predictions are meaningful.
\item \emph{Conditioning overfits the training distribution.} Mitigate by
keeping $g_\theta$ deliberately tiny and evaluating primarily on held-out
centres.
\item \emph{Gumbel selection adds gradient variance.} Mitigate by annealing $m$
downward from 44, so selection is introduced only after the weights are
informative.
\item \emph{The premise is wrong and per-sample weights collapse to a global
solution.} This is a real possibility. It is also a publishable negative result
about the structure of multi-scale supervision, and the ablation in \S4(iv) is
designed to detect it rather than hide it.
\end{itemize}

\textbf{Timeline (first 12 months).} Months 1--3: bounded operations, corrected
\code{concat}, reproduce LoMix at full budget as the reference point. Months
4--6: conditioning and calibration objective; in-distribution results on
Synapse/ACDC. Months 7--9: uncertainty-guided selection and the cost study;
cross-centre polyp protocol. Months 10--12: ablations, writing, submission
(MICCAI or ISBI, with the efficiency study a natural fit for a CVPR/ICCV
workshop).

\section{Fit with the group}

The proposal extends a lab paper along the lab's own axis --- efficiency without
accuracy loss --- and adds the trustworthiness dimension the group's stated
direction emphasises. It reuses the existing LoMix/EMCAD codebase rather than
starting elsewhere, and the first three months produce something useful to the
group regardless of whether the main hypothesis holds: a corrected operation set
and a full-budget reference run.

It also matches what I have been doing. \textbf{CLEAR-MoE} converts frozen
Vision Transformers into sparse mixture-of-experts models using a
calibration-driven layer-selection criterion --- the same ``use calibration to
decide where to spend capacity'' idea, applied to architecture rather than
supervision. \textbf{Predictive Self-Conditioning} conditions a LiDAR detector
on its own predictions and is evaluated on calibration precisely because
self-conditioning compounds error when confidence is wrong. CALM is the same
question in the lab's domain: when a model weights its own multi-scale
predictions, that weighting should depend on how much those predictions deserve
to be trusted for \emph{this} input.

\end{document}
